\section{Metodologia}

\subsection{Ferramentas e organização dos experimentos}

O pipeline foi implementado em Python 3.11.15 em três etapas: cálculo clássico das
energias de referência com PySCF 2.13.0 \cite{sun2018pyscf}; mapeamento para
qubits, gerenciamento de espaços ativos e execução do VQE com Qiskit Nature 0.7.2
\cite{qiskitnature2026}, Qiskit 2.3.0 \cite{javadi2024qiskit} e Qiskit Algorithms
0.4.0; e treinamento e validação dos modelos com scikit-learn 1.9.0
\cite{pedregosa2011scikit}. O código-fonte terá repositório público revelado após
a revisão duplo-anônima.

\subsection{Etapas e conjunto de dados}

\begin{figure}[ht]
  \centering
  \includegraphics[width=\textwidth]{pipeline.png}
  \caption{Etapas de construção do conjunto de dados e recomendação top-$k$. A
  energia FCI/CASCI é usada somente para produzir o rótulo de erro.}
  \label{fig:pipeline}
\end{figure}

Como ilustra a Figura~\ref{fig:pipeline}, para cada geometria calculamos uma
energia clássica de referência e executamos as 18 configurações da busca em
grade do VQE. Os erros obtidos formam os dados de treinamento. Depois de
treinado, o recomendador ordena os candidatos de uma geometria e indica apenas
os $k$ primeiros para execução.

Foram usadas H$_2$, LiH e BeH$_2$ na base STO-3G, com nove geometrias para cada
molécula. A Tabela~\ref{tab:distances} apresenta as distâncias avaliadas.

\begin{table}[ht]
\centering
\caption{Distâncias internucleares avaliadas, em \AA.}
\label{tab:distances}
\begin{tabular}{@{}lp{0.78\textwidth}@{}}
\toprule
Molécula & Distâncias \\
\midrule
H$_2$   & 0,3438; 0,5521; 0,7778; 0,9861; 1,2118; 1,4201; 1,6285; 1,8542; 2,0625 \\
LiH     & 0,8125; 1,2639; 1,7563; 2,2898; 2,7822; 3,3157; 3,8491; 4,3826; 4,8750 \\
BeH$_2$ & 0,5000; 0,9615; 1,2563; 1,5088; 1,7769; 2,1717; 2,6136; 3,1818; 3,7500 \\
\bottomrule
\end{tabular}
\end{table}

Para H$_2$, mantivemos o espaço completo de dois elétrons e dois orbitais
espaciais. Em LiH e BeH$_2$, os elétrons mais internos foram mantidos fixos para
reduzir o tamanho do problema. Com isso, usamos os espaços ativos CAS(2e,3o) e
CAS(4e,4o), respectivamente, em que ``e'' indica elétrons e ``o'', orbitais. A
transformação de Jordan--Wigner resultou em 4 qubits para H$_2$, 6 para LiH e 8
para BeH$_2$.

\begin{table}[ht]
\centering
\caption{Resumo do conjunto de dados. A última coluna indica execuções que
atingiram precisão química.}
\label{tab:dataset}
\begin{tabular}{lrrrr}
\toprule
Molécula & Qubits & Geometrias & Execuções & Precisão química \\
\midrule
H$_2$   & 4 & 9 & 162 & 71 (43,8\%) \\
LiH     & 6 & 9 & 162 & 50 (30,9\%) \\
BeH$_2$ & 8 & 9 & 162 &  9 ( 5,6\%) \\
\midrule
Total   & -- & 27 & 486 & 130 (26,7\%) \\
\bottomrule
\end{tabular}
\end{table}

A Tabela~\ref{tab:dataset} resume as 486 execuções. Em cada geometria, combinamos
os ansätze RealAmplitudes, EfficientSU2, UCCSD e PUCCSD com os otimizadores
COBYLA, SLSQP e SPSA. RealAmplitudes e EfficientSU2 usam uma ou duas repetições,
rotações padrão do Qiskit, emaranhamento linear e todos os qubits inicialmente
em zero. UCCSD e PUCCSD usam uma camada e partem do estado eletrônico de
referência gerado na preparação da molécula. Os parâmetros iniciais são
amostrados de $U[-\pi,\pi]$ com semente 137. Mantivemos as configurações padrão
dos otimizadores e limitamos cada execução a 100 iterações. As simulações usam
vetor de estado exato, sem amostragem (\textit{shots}) e sem ruído.

O modelo prevê o logaritmo do erro energético,
$y=\ln(1+|E_{\mathrm{VQE}}-E_{\mathrm{ref}}|_{\mathrm{kcal/mol}})$. Essa
transformação reduz o peso dos poucos erros muito altos durante o treinamento.
As entradas categóricas são molécula, ansatz e otimizador; as numéricas são
distância, número de repetições, qubits e termos da Hamiltoniana. Energias,
erros, tempos e avaliações do otimizador são excluídos porque só seriam
conhecidos após a execução e causariam vazamento de informação.

\subsection{Modelos, referências de comparação e validação}

Comparamos DummyRegressor, Ridge, RandomForestRegressor e MLPRegressor. Aplicamos
codificação \textit{one-hot}
às categorias e padronização às variáveis numéricas dentro de cada divisão. A
busca de hiperparâmetros considera MLPs com camadas $(32,16)$ ou $(64,32)$ e
$\alpha\in\{10^{-3},10^{-2}\}$; Ridge com
$\alpha\in\{0{,}1,1,10\}$; e Random Forest com 300 árvores, profundidade livre
ou 8 e folhas mínimas de 1 ou 2.

Na avaliação principal, os dados são separados cinco vezes, sempre agrupados
por molécula e distância. Assim, as 18 configurações de uma geometria ficam
juntas no treino ou no teste. Em cada separação, os hiperparâmetros são escolhidos
com outras três divisões, também agrupadas. O recomendador é comparado com a
seleção aleatória, a melhor configuração geral e a melhor configuração de cada
molécula, calculadas apenas com os dados de treino. Também mostramos o melhor
resultado possível quando toda a busca é conhecida. Para entender a importância
do nome da molécula, repetimos a avaliação sem essa informação. Por fim, outro
teste retira uma molécula inteira do treino.

MAE e RMSE são calculados em kcal/mol após a transformação inversa; MAE e
$R^2$ também são calculados na escala logarítmica de $y$. Para top-1, top-3 e
top-5, medimos quantas
geometrias atingem precisão química, o menor erro entre os candidatos escolhidos
e a diferença para o melhor candidato disponível. A fração $(18-k)/18$ indica
somente quantos candidatos deixam de ser testados, não a economia de tempo ou de
custo quântico. Os intervalos de confiança usam mil reamostragens
(\textit{bootstrap}). Na comparação entre MLP e as estratégias simples, as
reamostragens preservam os pares de resultados de cada geometria. Também fazemos
uma análise mais conservadora, reamostrando as três moléculas como grupos.